\scanpage{266}
\mspage{- 84 -}\hfill IX

\uline{Proposition} 7.1.5. \uline{Avec les notations précédentes pour} $S,A,\ \hat B$, \uline{supposons que} $S$ \uline{soit} normal. \uline{Alors} $\hat B$ \uline{est algébrisable}, \uline{et le schéma abélien} $B$ \uline{sur} $S$ \uline{dont il provient est projectif sur} $S$.

On peut évidemment supposer $A=A^{\unsure{o|0|\circ}}$, et $S$ intègre. D'après M. RAYNAUD [23 XI 1.13], l'hypothèse que $S$ est normal implique que $A$ est quasi-projectif sur $S$. Soit donc $\underline L$ un Module inversible sur $A$ qui est relativement ample, et formons le Module birigidifié $\delta(\underline L)=\pi^*(\underline L)\operatorname{pr}_1^*(\underline L)^{-1}\operatorname{pr}_2^*(\underline L)^{-1}$ sur $A\times_S A$ comme dans VIII 4.12. Comme les fibres de $A_{\unsure{o|0}}$ sont de rang unipotent nul, il en est de même de celles de $A$ (résulte facilement de 3.1), donc en vertu de VIII 7.5, $\delta(\underline L)$ définit une biextension $W$ de $(A,A)$, évidemment symétrique, qui en vertu de VIII 4.13 (i.e. essentiellement en vertu de RAYNAUD [23 XI 1.11]) satisfait à la condition d'amplitude voulue en les points de $S_o$.

\uline{Remarque} 7.1.6. On peut prouver, sans hypothèse de normalité sur $S$, que $\hat B$ est toujours algébrisable, en se ramenant au cas normal comme dans 13.5.3 a).

7.2. Notons maintenant le

\uline{Lemme} 7.2.1. \uline{Soient} $S$ \uline{le spectre d'un anneau noethérien} $V$, \uline{séparé et complet pour la topologie définie par un idéal} $\underline m$ \uline{de} $V$, $C$ \uline{la catégorie des schémas en groupes} $A$ \uline{sur} $S$, \uline{lisses de type fini sur} $S$, \uline{tels que} $A/A^o$ \uline{soit un schéma en groupes étale fini sur} $S$ \uline{et que} $A^o$ \uline{soit une extension d'un schéma abélien} $B$ \uline{par un tore isotrivial} $T$ (SGA 3 IX 1.1),

\handwritten{\unsure{402|407}}
\mspage{402}
% archive sheet number: 265
